\section{Frozen decompositions and cycle quality}
\label{sec:backend-quality}

We now construct the flat-cycle backend of
Lemma~\ref{lem:flat-cycle}.  This section proves its quality guarantee:
one recursive level reduces an arbitrary current circulation to local
correction cycles and a circulation in one smaller graph.  The construction
uses verified covers of a frozen graph, so its conclusion holds for every
subsequent state of an epoch, including states chosen adaptively.
Section~\ref{sec:backend-updates} supplies the path assignments, maintains
the smaller graph, and publishes a forest in which the returned walks have
the required representation and recourse.

\subsection{Parameters, scores, and the bottom level}

Continue with the parameters of Lemma~\ref{lem:routing-helper}:
\begin{equation}
\label{back:helper-parameters}
 b=M^{1/a^2},\qquad
 \Lambda=2^{1000a^2}(1+u),\qquad
 P_*=\Lambda^{400a^2},\qquad M=2^u.
\end{equation}
Here $u$ is a positive multiple of $2a^2$ and
$S^{1/a}\le M\le S^2$.  All powers in the following display are
integers:
\begin{equation}
\label{back:parameters}
\begin{gathered}
 d=2a-4,\qquad k=M^{1/a},\qquad
 n_i=M^{1-i/(2a)},\qquad G_i=1200a+2i
       \quad(0\le i\le d),\\
 r=(1210a+3)a^2,\qquad \Gamma=K_h=\Lambda^2,\\
 A_*=64r\Gamma K_h b,\qquad
 p_*=(10P_*)^{a^2+1},\qquad Z_*=40(p_*+1).
\end{gathered}
\end{equation}
In addition to the helper's prerequisites, impose the deterministic checks
\begin{equation}
\label{back:guards}
                 P_*\le b,\qquad 20Z_*A_*\le b^2.
\end{equation}
For each fixed $a$, these inequalities hold for all sufficiently large
$M$: every factor other than the displayed powers of $M$ is a fixed
power of $1+u$.  A failed prerequisite suspends the parameter branch.

Level $i$ accepts an undirected multigraph on fixed vertex slots, of total
size at most $n_i$, counting slots and edge copies.  Its lengths are
positive integers and its signed gradients are integers, with
\[
                   1\le\ell_e\le M^{G_i},\qquad
                   |\gamma_e|\le M^{G_i}.
\]
Each edge has an identifier and a reference orientation.  The degree
bound needed for the running time is a fixed-parameter logarithmic
power; Section~\ref{sec:backend-updates} proves that this bound persists
under recursion.  The outer input has degree at most three.
Write $\OPT_i$ for the maximum circulation ratio on the current
level-$i$ graph, with the same zero convention as in
Lemma~\ref{lem:flat-cycle}.

A \emph{certified score} is a nonnegative rational lower bound on the
absolute gradient divided by the positive raw length of a returned
traversal list.  A zero score may accompany no candidate.  Raw length
counts every actual traversal with multiplicity; in the constructions
below, assigning length one to an artificial null link only enlarges the
length used for an upper bound.  When a vector is used solely in an
existential comparison, its length norm is instead the sum of length
times the absolute \emph{net} coefficient on each edge.  We explicitly
distinguish these two uses throughout.

\begin{lemma}[Bottom-level search]
\label{back:base}
At level $d$, a query can return individual oriented edges forming a
circulation with certified score at least $\OPT_d/2$, whenever
$\OPT_d>0$.  Initialization and each graph operation or query cost
$M^{O(1/a)+o_a(1)}$, and the answer has at most $n_d=M^{2/a}$ pieces.
\end{lemma}

\begin{proof}
Replace each undirected edge by its two possible traversals.  For each
dyadic threshold $\rho=2^j$ with
\[
              -(G_d+2)u\le j\le (G_d+1)u,
\]
give these two arcs costs $\rho\ell_e-\gamma_e$ and
$\rho\ell_e+\gamma_e$, respectively.  For every start vertex, tabulate
minimum costs of walks with each number of arcs from zero through the
number of vertices, retaining predecessor pointers.  This detects a
negative closed walk whenever one exists, because a negative closed walk
contains a negative simple directed cycle.  Keep a witness for the largest
successful threshold, and return its exact absolute-gradient to raw-length
ratio.

To justify the approximation, negate a nonzero signed circulation if
necessary to make its gradient nonnegative.  Orient its used edges along
their coefficient signs and decompose the resulting nonnegative directed
circulation into simple directed cycles.  Gradient and length add in this
decomposition, so one cycle has ratio at least that of the circulation.
Conversely, any positive-ratio directed cycle projects to a nonzero signed
circulation; its projected length norm is at most its raw length.  It
follows that $\OPT_d$ is attained by a directed cycle.  Integer gradients
and positive integer lengths give
\[
        \frac{1}{n_d M^{G_d}}\le\OPT_d\le M^{G_d}
                  \qquad\text{if }\OPT_d>0.
\]
The tested thresholds cover this interval with room on both sides.  The
largest successful threshold is at least $\OPT_d/2$, and its negative
walk has still larger ratio.  The witness has at most the number of
vertices in the graph.

There are $O_a(1+u)$ thresholds.  The elementary tables take at most
$O(n_d^3)$ arithmetic operations per threshold, since both the vertex
and arc counts are $O(n_d)$.  All costs have $O_a(1+u)$ bits after a
common dyadic scaling.  This proves the stated bound, including charged
multiword arithmetic, without a forest at the bottom level.
\end{proof}

For $i<d$, use epochs of
\begin{equation}
\label{back:epoch}
                         h_i=n_i/k
\end{equation}
single input edits.  A price replacement is a deletion followed by an
insertion.  At the beginning of an epoch, freeze the current nonloop
graph $\bar G$, including its prices and edge identities.  An old edge
keeps these data until its death.  Every later insertion is new, even
when its endpoints and prices coincide with those of a dead edge.
A heap on all current loops supplies exact single-loop candidates.
Loops can otherwise be omitted: their coordinates are independently
circulations, and the optimum on a graph is the maximum of its loop
optimum and its nonloop optimum.

Fix a current nonloop circulation $\Delta$, and write
$\Delta=\Delta_s+\Delta_n$ on the surviving frozen and new edges.
The surviving part need not be a circulation: its imbalance is the
negative of that of $\Delta_n$, so it is supported at new-edge endpoints.
We will express $\Delta$ as a linear combination of current circulation
candidates plus the projection of one circulation in a smaller graph.
The coefficient-weighted lengths and the child vector's length norm must
be bounded by a controlled multiple of $\|\ell\Delta\|_1$.
Each candidate, however, needs a bound on its own raw traversal length,
because that is the cost charged when the engine uses it.

\subsection{Verified covers and a finite family of forests}

The covers below adapt the shortest-path partitions with random shifts
of Miller, Peng, and Xu
\cite[Algorithm~2 and Lemmas~4.1 and~4.4]{RandomShifts}.
We prove the required guarantees for weighted distances, geometric
delays, verified padding, and capped retries.

For each layer $j=1,\ldots,r$, let $R_j=b^j$.

\begin{lemma}[Verified padded covers]
\label{back:covers}
For every layer $j$ one can construct $K_h$ partitions of the frozen
vertices into cells, with a shortest-path tree in each cell, such that:
\begin{enumerate}[label=(\roman*)]
\item every cell tree has radius at most $\Gamma R_j$ about its center;
\item for every vertex, at least one of the partitions has a cell
containing its entire closed $R_j$-ball in $\bar G$.
\end{enumerate}
The properties are verified before use.  The construction takes
$n_iM^{o_a(1)}$ time on every path for fixed reliability parameter
$\sigma$, with detected-failure probability at most
$2^{-\sigma}S^{-800}$ conditional on any valid frozen input.  It uses
integer operations and fresh fair random bits.
\end{lemma}

\begin{proof}
Independently for each partition, give every vertex $c$ an integer delay
$X_c$ with $\Pr[X_c\ge t]=2^{-t}$ for nonnegative integers $t$.
At $v$, choose the center maximizing
\begin{equation}
\label{back:shifted-score}
                        R_jX_c-\dist_{\bar G}(c,v),
\end{equation}
breaking ties by center identifier.  Reject the attempt if any delay
exceeds $\Gamma$.  A vertex's own center has nonnegative score, so its
winning center is at distance at most $\Gamma R_j$.  Dominance of a
center persists along a shortest path towards it: the triangle
inequality preserves every score comparison, including the common
identifier tie rule.  Hence a shortest-path tree to the center lies
inside the winning cell.

For each vertex also compute the runner-up among distinct centers.
Verify that at each layer every vertex has, in some partition, winning
margin greater than $2R_j$.  The difference between the scores of two
centers changes by at most twice the distance moved.  Thus this check
implies containment of the whole closed $R_j$-ball.  A runner-up in
another connected component has score minus infinity and causes no
difficulty.

For unlimited geometric draws, fix a vertex and all delays except that
of a proposed winner $c$.  The event that $c$ wins is a lower integer
threshold on $X_c$, truncated below at zero.  Exceeding that threshold
by three raises its score by $3R_j$ and gives margin greater than
$2R_j$.  The geometric tail assigns probability at least $1/8$ of the
winning-event probability to this overshoot.  Summing over the disjoint
winning events gives margin probability at least $1/8$.  Independence
of the partitions and a union bound therefore bound the probability
that the delay or coverage checks fail by
\begin{equation}
\label{back:cover-failure}
       rK_hn_i2^{-\Gamma}+rn_i(1-1/8)^{K_h}<\frac12.
\end{equation}
The displayed inequality follows from the chosen scales in the checked
parameter range.  Repeat with fresh draws at most
$2(\sigma+1000\log_2 S+10)$ times.  Exhaustion is a detected failure;
its probability is at most $2^{-\sigma}S^{-800}$.

Generate a geometric delay by fair bits, stopping and rejecting when
its cap is exceeded.  Shifted distances and the best two distinct
center labels are computed by multi-source Dijkstra with initial
offsets $-R_jX_c$, finalizing at most two distinct origins at each
vertex.  An origin among the best two at a vertex has prefixes among
the best two along a shortest path; otherwise two strictly preferred
origins at a prefix would remain preferred at the endpoint.  Thus this
truncation loses neither of the required origins nor a winning tree
pointer.  Initial offsets can be negative because edge lengths are
positive.  Each partition needs only near-linear integer graph work.
The layer, partition, delay, and repetition counts are fixed powers of
$1+u$ for fixed $a,\sigma$, and their integer data have
fixed-parameter logarithmic bit length.  This proves the every-path
cost bound.
\end{proof}

Let $H_j$ be the disjoint union of copies of all the cell trees in the
$K_h$ partitions at layer $j$.  A tree edge maps to its frozen actual
edge, with the appropriate endpoint correspondence and reference
orientation.  Introduce additional vertices $V_j(v)$ for every frozen
slot $v$ and $0\le j\le r$.  The layered system has two kinds of
artificial links:
\begin{itemize}
\item an OUT link from $V_{j-1}(v)$ to the copy of $v$ in each partition
of $H_j$;
\item an IN link from every cell root $c$ to $V_j(c)$.
\end{itemize}
Each artificial link maps to a null traversal at one physical slot,
and is assigned length one and gradient zero for the construction.

The following finite family adapts the bit-selection construction of
Chen et al.\ \cite[Definition~4.30 and Lemma~4.31]{IncrementalGraphs}.

\begin{lemma}[Finite family of routing forests]
\label{back:forest-family}
There is a family of at most $(K_h^2(u+1))^{O(r)}$ forests $T^0$ with
the following property.  At each layer, any two distinct lower slots
can be prescribed arbitrary OUT options, and some family member
realizes all these prescriptions simultaneously across the layers.
Each member contains separate copies of all cell trees and IN links,
and exactly one OUT link per lower slot.
\end{lemma}

\begin{proof}
At a given layer choose a bit position of the slot index and two
partition options, using the chosen option as a function of that bit.
Include constant functions as well.  Two distinct indices differ in
some bit, so any two prescribed options can be realized.  Choose these
functions independently at the $r$ layers and include every choice in
the family.  There are at most $2K_h^2(u+1)$ choices per layer, giving
the claimed bound.

Orient each cell tree towards its root, and the links towards higher
layers.  Every vertex has outdegree at most one and there is no directed
cycle.  Such an underlying undirected graph is a forest: an undirected
cycle would, by its edge count and outdegree bound, force a directed
cycle.  This also proves the asserted structure of each $T^0$.
\end{proof}

The chords of each $T^0$ are the frozen input nonloops on its $V_0$
vertices, and a null pair chord between every two distinct OUT-option
targets of the same lower vertex.  The latter has length one and
gradient zero, whether or not either OUT option was selected in $T^0$.
Retain only chords whose endpoints are connected in that forest.
For such a chord $f$, let $C_f^0$ be its signed fundamental cycle,
traversing $f$ forward and the tree path backwards.  Let $w_f$ be its
full length, including artificial links, and fix the bucket
\begin{equation}
\label{back:bucket}
                         W_f=2^{\lceil\log_2 w_f\rceil}.
\end{equation}
Its value is found by integer bit-length operations, so logarithms in
this notation require no real arithmetic.

Whenever a frozen actual edge dies, delete every tree edge mapping to
it, both in the $T^0$ copies and in the separate raw forests $H_j$.
These deletions are called \emph{holes}.  For a surviving chord $f$,
write $C_f^{\mathrm{clip}}$ for the chord together with the surviving
signed tree traversals of $C_f^0$.  Null chords always survive.
The clipped vector need not be a circulation.

\subsection{Fixed marginals and the clipped decomposition}

Let $\pi$ denote projection onto the frozen actual edge coordinates,
ignoring null links and adding signed coefficients of copies.  We
first define a distribution of an upward route for every vertex.
These distributions are used only in the proof, and are not sampled
or evaluated by the algorithm.

The normalized boundary-distance weights and retention of each starting
vertex's route law adapt the origin-aware routing of Rozho\v{n} et al.\
\cite[Section~4]{MinorAggregateRouting} and its dynamic development in
Chen et al.\ \cite[Section~4.3]{IncrementalGraphs}.

To account for the imbalance of $\Delta_s$, we will compare a surviving
edge with the difference of two upward routes, one from each endpoint.
The same vertex must have the same route law when it appears at the
endpoint of different edges.  With this fixed marginal, the route
contributions cancel wherever $\Delta_s$ is conserved; only its imbalance
at new-edge endpoints remains.  Thus the route law will be indexed by
the starting vertex, even when the route has reached a different
representative in a higher layer.

For a cell $C$ at layer $j$, including its partition identity, define
\begin{equation}
\label{back:cell-weight}
 a_{j,C}(v)=
 \max\!\left\{0,
   \min\{R_j,\dist_{\bar G}(v,V(\bar G)\setminus C)\}
                       -\frac{R_j}{2}\right\}.
\end{equation}
Distance to the empty set, or to vertices in other components only,
is infinity.  Normalize these weights to a distribution
$\mathcal L_j(v)$ on all cells at that layer.  A verified covering
cell gives a weight of at least $R_j/2$, so the normalizer is at least
$R_j/2$.  A positive-weight cell contains every vertex at distance at
most $R_j/2$ from $v$.

\begin{lemma}[Upward route laws]
\label{back:path-law}
The independent layer choices $\mathcal L_j(v)$ define an upward path
$P(v)$ in the layered system with all OUT options.  For every frozen
edge $e=vw$, its two path laws admit a coupling such that the closed
walk $e+P(w)-P(v)$, after canceling a common final suffix, has expected
full length at most $A_*\ell_e$.  The marginal law of $P(v)$ is the
same for every incident edge used in this assertion.
\end{lemma}

\begin{proof}
The functions in \eqref{back:cell-weight} are one-Lipschitz on a
connected component.  At two vertices, at most $2K_h$ cell weights
can be nonzero, since a positive weight requires membership in the
cell.  Comparing normalized weights along $e=vw$ yields
\begin{equation}
\label{back:tv}
              \TV(\mathcal L_j(v),\mathcal L_j(w))
                        \le\frac{4K_h\ell_e}{R_j}.
\end{equation}
Indeed the sum of absolute unnormalized changes is at most
$2K_h\ell_e$, and total variation of the normalized laws is at most
that sum divided by the smaller normalizer.

Starting at $V_0(v)$, use the chosen cell at layer one, then its tree
path to its center and the IN link to the corresponding $V_1$ vertex.
Continue likewise.  At layer $j\ge2$ the previous representative is
the center of a cell containing $v$, so its distance from $v$ is at
most $\Gamma R_{j-1}$.  The second check in \eqref{back:guards} implies
$b>2\Gamma$, and therefore
\[
                  \Gamma R_{j-1}<R_j/2.
\]
Every positive next-layer choice contains that representative, so its
OUT option exists.  This defines $P(v)$ using independent choices
across layers and a marginal law fixed by $v$ alone.

For the endpoints of $e$, couple their cell choices maximally at each
layer, independently between layers.  At the top,
$R_r/2$ exceeds every finite component diameter: such a diameter is
at most $n_iM^{G_i}$, whereas $R_r=M^{1210a+3}$.
A top-layer cell with positive weight at a vertex must consequently
contain its whole connected component.  All such cells have weight
$R_r/2$ at every vertex of that component.  The two endpoint laws
therefore coincide at the last layer and can be coupled to agree,
including the partition identity.

Retain the two arms through one layer after their last cell
disagreement; if there is no disagreement, retain layer one.  From
then on their representative and choices coincide, so their suffixes
are identical and cancel.  Two arms at layer $j$, including their
OUT and IN links, cost at most $4\Gamma R_j$.  For $j\ge2$ this layer
is retained only if a disagreement occurred at some layer $l\ge j-1$.
Equation~\eqref{back:tv} and a geometric sum bound that probability by
\[
       \sum_{l\ge j-1}\frac{4K_h\ell_e}{R_l}
                         \le\frac{8K_h\ell_e}{R_{j-1}}.
\]
Including the original edge and the first layer, the expected length
is consequently at most
\[
 \ell_e+4\Gamma b+
       \sum_{j=2}^r32\Gamma K_hb\ell_e
                   \le64r\Gamma K_hb\ell_e=A_*\ell_e.
\]
The coupling changes none of the prescribed per-vertex marginals.
\end{proof}

The next decomposition uses the pair-chord construction of
\cite[Definitions~4.28 and~4.30 and Lemma~4.31]{IncrementalGraphs}.

Each coupled walk decomposes into fundamental cycles from the finite
family.  To see this, split the two trimmed arms at their common
representative vertices $V_j$.  Before the first common vertex, the
two lower slots at every layer are distinct.  Lemma~\ref{back:forest-family}
therefore supplies a forest containing the paths that close the input
chord $e$.  Between successive common vertices, equal first OUT
options give identical tree paths to the next layer and cancel.  If
those first options differ, replace the two null OUT traversals by
the corresponding null pair chord.  For the remaining arms in that
piece, all simultaneous lower-slot prescriptions concern distinct
slots; another family member realizes them.  The pair-chord endpoints
are connected there, as the arms meet at the next common representative.
Simplifying each tree walk preserves its signed tree vector and can
only decrease length.  Replacing two unit null links by one unit pair
chord also cannot increase length.  Thus the projected coupled walk
is a sum of projected oriented fundamental cycles whose full lengths
sum to at most its original length.

\begin{lemma}[Frozen decomposition with holes]
\label{back:frozen-decomposition}
Fix any state of the epoch and any circulation $\Delta$ on its current
nonloops.  Write $\Delta=\Delta_s+\Delta_n$ for its surviving frozen
and new coordinates.  There are real coefficients $\alpha_f$ on live
fundamental chords and a vector $z$ on the unduplicated raw forests
$H_j$ such that, before restricting the frozen coordinates,
\begin{align}
 \sum_f\alpha_f\pi C_f^0&=\Delta_s+\pi z,
                                      \label{back:frozen-identity}\\
 \sum_f|\alpha_f|w_f&\le A_*\|\ell\Delta_s\|_1,
                                      \label{back:coefficient-bound}\\
 \|z\|_{\mathrm{length}}&\le A_*\|\ell\Delta_s\|_1.
                                      \label{back:raw-residual-bound}
\end{align}
Moreover $z$ is the raw-tree part of
\begin{equation}
\label{back:marginal-cancellation}
 \sum_v\bigl(\operatorname{in}_{\Delta_s}(v)
                  -\operatorname{out}_{\Delta_s}(v)\bigr)\E[P(v)].
\end{equation}
Consequently this description uses only upward paths starting at slots
incident to new edges.  If $z_L$ is the restriction of $z$ to live raw
forest edges, then
\begin{equation}
\label{back:clipped-identity}
             \sum_f\alpha_f\pi C_f^{\mathrm{clip}}-\pi z_L
                                  =\Delta_s.
\end{equation}
\end{lemma}

\begin{proof}
Apply the preceding coupled-walk decomposition separately to each
surviving frozen edge, weighted by its signed coefficient in
$\Delta_s$, and take expectations.  There are finitely many choices
in these conceptual distributions.  Collecting coefficients of the
oriented family chords gives $\alpha_f$; collecting the signed route
parts on the raw copies gives $z$.  Only surviving input edges and
null pair chords are used as chords, so every indexed chord is live.
The triangle inequality and Lemma~\ref{back:path-law} give
\eqref{back:coefficient-bound} and \eqref{back:raw-residual-bound},
as well as the exact identity \eqref{back:frozen-identity}.

Although the couplings may differ between incident edges, the law of
each endpoint path is fixed.  Summing the route term
$P(w)-P(v)$ over oriented edges therefore gives
\eqref{back:marginal-cancellation}.  Cancellation of identical final
suffixes does not change this signed vector.  Because $\Delta$ is a
circulation, the displayed imbalance of $\Delta_s$ is the negative
of that of $\Delta_n$ and vanishes away from new-edge endpoints.

Finally delete every coordinate of every tree copy mapped to a dead
frozen edge.  Projection after this deletion is the same as restriction
of the projected physical vector.  Restricting
\eqref{back:frozen-identity} therefore proves
\eqref{back:clipped-identity}, even when coefficients on a dead edge
had canceled between different copies.  No circulation property is
asserted for an individual clipped cycle.
\end{proof}

The identity~\eqref{back:clipped-identity} accounts for every surviving
frozen coordinate.  Its fundamental-cycle terms need not be current
circulations, because holes may have opened them.
We next extract current circulation candidates using short
routed corrections and compress the remaining forest vectors.
The description of $z$ uses only upward routes starting at new-edge
endpoints.  The endpoints of their cell-tree legs, together with hole
endpoints, will supply the compression terminals for this raw-forest
remainder.

\subsection{Portal legs and local correction cycles}

\begin{lemma}[Portal grid]
\label{back:portals}
A forest admits a set of at most its number of vertices divided by $k$
portals such that every component remaining after portal removal has
fewer than $k$ vertices.  The set is constructed by one postorder
traversal.
\end{lemma}

\begin{proof}
Root each tree arbitrarily.  Accumulate the number of still-connected
unmarked vertices in the subtrees of a vertex, including the vertex
itself.  If the accumulated count is at least $k$, mark that vertex as
a portal and pass count zero to its parent.  Every mark charges a
disjoint group of at least $k$ vertices.  In any component off the
marks, the residual count at its highest vertex is exactly its size
and is below $k$.
\end{proof}

Place such a grid on every $T^0$ and separately on each raw forest.
Call the components off the grid \emph{pockets}.  Grids remain
compression terminals for the whole epoch.  Classify and assign each
live fundamental chord separately in its forest and bucket $W$.
\begin{enumerate}[label=(\roman*)]
\item If its full frozen tree path touches the grid, call the chord
\emph{regular}.  At each incidence choose the nearest portal to its
endpoint in the full frozen tree metric, breaking ties by a fixed
portal order.  Use the leg towards that portal, truncated at the first
hole and ending on the incidence side of that hole.  The full path
contains a portal at distance at most $w_f\le W$ from each endpoint,
so each such leg has length at most $W$.
\item Otherwise the full tree path lies in one pocket.  Root that
pocket and measure weighted depth.  For this bucket take connected
clusters in half-open depth bands of width $4W$, using both the
unshifted bands and their translate by $2W$.  The depth range of the
chord path has width at most its path length, hence at most $W$; one
of the two band systems contains that whole range.  Assign the chord
to its containing connected cluster there.  Its tree diameter is at
most $8W$, since the least common ancestor of two vertices in a
connected band cluster is in that cluster.

Until any interior edge of that cluster is cut, use the intact
fundamental cycle directly.  At its first such cut, activate all
surviving assigned chords.  Give each of the two resulting sides a
class with target its endpoint of the cut.  A later cut splits the
corresponding live fragment into two fresh classes, with targets the
new cut endpoints.  All resulting legs stay inside the original
cluster and have length at most $8W$.
\end{enumerate}

A \emph{class} is a logical vertex whose incidences have a common
target tree node.  Regular incidences are initially grouped by their
target portal; holes subsequently split those classes.  Splits always
create fresh child labels, including for an empty side.  The exact
class maintenance is given in Lemma~\ref{update:classes}.  A chord
using the leg scheme is called \emph{active}.

For every active $f$, the construction in
Section~\ref{sec:backend-updates} supplies a logical path $\Pi_f$
between its endpoint classes, using selected active chords of the
same forest and bucket and containing at most $p_*$ hops.  The
assignments may use an empty path for a logical loop.  If $f$ is
oriented from its tail to its head, let $L_f$ traverse from its tail
target to its tail, then across $f$, and then from its head to its head
target.  Such walks concatenate, also in reverse, to realize logical
paths; write $L_{\Pi_f}$ for the resulting walk.  Every logical hop
has full length at most
\begin{equation}
\label{back:logical-hop}
                         8W+W+8W=17W.
\end{equation}

The local correction candidate is
\begin{equation}
\label{back:correction}
 B_f=\begin{cases}
       C_f^0,&\text{while the chord is used directly},\\
       L_f-L_{\Pi_f},&\text{for an active chord}.
     \end{cases}
\end{equation}
Here a negative walk is traversed backwards; this formula specifies a
traversal list, not a cancellation rule for its denominator.  Its
projection is a circulation, and its raw actual length is at most
\begin{equation}
\label{back:correction-length}
             17(p_*+1)W_f\le Z_*W_f.
\end{equation}
The direct case satisfies the same bound because $w_f\le W_f$.
Hence
\begin{equation}
\label{back:local-score}
                     \frac{|\gamma^{\mathsf T}\pi B_f|}{Z_*W_f}
\end{equation}
is a certified score.  A positive numerator ensures positive actual
raw length: null traversals have zero gradient.  The list consists of
actual individual chord edges and live forest paths; null chord
traversals can be omitted.  Its numerator is computed from integer
path-gradient sums, with the fixed positive denominator above.

\subsection{One joint child graph}
\label{back:child}

The smaller instance is one graph $J$ incorporating all family
forests and the separate raw forests.  In each live component take
the Steiner subtree of its terminals, retain the terminals and
branching vertices, and suppress degree-two nonterminals.  A
one-terminal component contributes its terminal alone; an empty
terminal set contributes nothing.  The terminal set includes:
\begin{enumerate}[label=(\roman*)]
\item every initial portal and both endpoints of every hole;
\item every class target and the \emph{original tree-node} endpoints
of every selected chord;
\item in the raw forests, for every slot ever incident to a new edge
in this epoch, both forest endpoints of every possible cell-tree leg
in its frozen upward route.
\end{enumerate}
All promotions may persist to the end of the epoch.  For item (iii),
enumerating all sequences of partition options suffices, even if some
enumerated routes have zero probability in the conceptual law.
Only their endpoints are needed; the cell-tree paths need not be
enumerated.  There are $K_h^{O(r)}$ such endpoints per incident slot,
a fixed-parameter logarithmic factor.

A compressed segment receives the sum of its constituent lengths
and the oriented sum of their gradients.  Include each currently
selected chord on its original tree-node endpoints, with its actual
or null prices; earlier selected chords may also remain while live.
For every needed physical slot introduce a hub, and join each
compressed vertex to the hub of its physical image by a unit-length,
zero-gradient null edge.  Include all current new nonloops directly
between their hubs.  Unused child slots remain isolated.

Every child edge expands into a live parent walk between the physical
images of its endpoints, with exactly its gradient and with raw
actual length at most its child length.  We also write $\pi$ for the
induced projection from child coordinates to current parent coordinates;
on surviving frozen edges it agrees with the earlier projection.
A compressed segment has at
most $k+1$ constituent edges: it has no internal portal and its
internal nodes lie in one pocket.  Consequently its length and
absolute gradient are at most
\[
                  (k+1)M^{G_i}\le M^{G_i+2}=M^{G_{i+1}}.
\]
The same bound holds for the other child edges.  The number of child
slots and edits, and their degree bounds, are proved in
Section~\ref{sec:backend-updates}; only the stated graph and maps are
needed for the following comparison.

\begin{lemma}[Compression of a forest vector]
\label{back:compression}
If a vector on a forest has divergence only at a specified terminal
set, it vanishes off the terminal Steiner subtrees and has constant
oriented coefficient along each suppressed segment.  Compression
therefore preserves its gradient and its length norm exactly.
\end{lemma}

\begin{proof}
Deleting an edge outside a terminal Steiner subtree leaves a side
containing no terminal.  The sum of divergences on that side is zero,
so the edge coefficient is zero.  At an internal degree-two
nonterminal of a remaining segment, zero divergence forces the two
consistently oriented coefficients to agree.  Replacing that segment
by an edge with summed length and signed gradient consequently
preserves both its gradient contribution and its length norm.
\end{proof}

For an active $f$, subtracting its correction from its clipped vector
leaves the clipped tree path minus the two legs of $L_f$, together
with $L_{\Pi_f}$.  The first part is a live forest vector with
divergence only at hole endpoints and the two targets.  The original
chord-endpoint divergences cancel, also when a hole is adjacent to
one of those endpoints.  All remaining divergence points are
terminals.  The second part uses only selected chords and live paths
between terminals.  By Lemma~\ref{back:compression} there is a child
vector $J_f$, not necessarily a circulation, satisfying
\begin{equation}
\label{back:remainder}
 \pi J_f=\pi(C_f^{\mathrm{clip}}-B_f),\qquad
                    \|\ell^J J_f\|_1\le Z_*W_f.
\end{equation}
For the length bound, the clipped tree part costs at most $W_f$,
the two removed legs at most $16W_f$, and the routed selected path
at most $17p_*W_f$.  Their sum is at most
$17(p_*+1)W_f$.  In the direct case the relevant cluster path is
intact, so $C_f^{\mathrm{clip}}=B_f$ and we take $J_f=0$.

Likewise, the vector $z_L$ in
Lemma~\ref{back:frozen-decomposition} compresses to a raw-forest child
vector $\widehat z_L$ with its length norm preserved.  Indeed its
description \eqref{back:marginal-cancellation} is a combination of
cell-tree legs whose endpoints were included in terminal item (iii).
Clipping creates additional divergence only at hole endpoints.
It is this \emph{net forest vector} that is compressed; no assertion
is made that its individual conceptual routes have disjoint support.

Figure~\ref{fig:frozen-compression} illustrates the compression of one
net forest contribution to the joint child graph.
\begin{figure}[tbp]
  \centering
  \begin{tikzpicture}[
    font=\small, >=Stealth,
    terminal/.style={circle,draw,fill=white,minimum size=5.5mm,inner sep=0pt},
    ordinary/.style={circle,fill=black,minimum size=2mm,inner sep=0pt},
    flow/.style={->,thick,blue!65!black},
    labelbox/.style={fill=white,inner sep=2pt}
  ]
    \node[anchor=west] at (0,1) {Live forest after one hole};
    \node[terminal,label=above:$t_-$] (ta) at (0,0) {};
    \node[ordinary] (va) at (1.5,0) {};
    \node[terminal,label=above:$h_1$] (ha) at (3,0) {};
    \node[terminal,label=above:$h_2$] (hb) at (6,0) {};
    \node[ordinary] (vb) at (7.5,0) {};
    \node[terminal,label=above:$t_+$] (tb) at (9,0) {};
    \draw[flow] (ta) -- node[above] {$\xi_1$} (va);
    \draw[flow] (va) -- node[above] {$\xi_1$} (ha);
    \draw[flow] (hb) -- node[above] {$\xi_2$} (vb);
    \draw[flow] (vb) -- node[above] {$\xi_2$} (tb);
    \draw[densely dashed,black!45] (ha) -- (hb);
    \draw[black!65,thick] (4.35,-.16) -- (4.65,.16);
    \draw[black!65,thick] (4.35,.16) -- (4.65,-.16);
    \node[below] at (4.5,-.25) {deleted edge};
    \node[ordinary] (leaf) at (1.5,-1.3) {};
    \draw[black!35] (va) -- node[right,text=black!60] {$0$} (leaf);
    \node[anchor=north,text=black!60,align=center] at (1.5,-1.48)
      {no terminal\\on this branch};
    \node[align=left,anchor=west,text width=3.2cm] at (9.75,-.5)
      {Open circles\\are terminals.\\[3pt]
       Both hole endpoints are retained.};

    \draw[->,semithick] (4.5,-1.1) --
      node[right,align=left,text width=4.5cm]
        {Steiner restriction and\\degree-two suppression} (4.5,-2.45);

    \node[anchor=west] at (0,-2.95) {Compressed contribution in $J$};
    \node[terminal,label=below:$t_-$] (cta) at (0,-3.7) {};
    \node[terminal,label=below:$h_1$] (cha) at (3,-3.7) {};
    \node[terminal,label=below:$h_2$] (chb) at (6,-3.7) {};
    \node[terminal,label=below:$t_+$] (ctb) at (9,-3.7) {};
    \draw[flow] (cta) -- node[above] {$\xi_1$} (cha);
    \draw[flow] (chb) -- node[above] {$\xi_2$} (ctb);
    \node[align=left,anchor=west,text width=3.2cm] at (9.75,-3.7)
      {Each segment stores summed length and oriented gradient.};
  \end{tikzpicture}
  \caption{Compression preserves a net tree flow whose divergence is
    supported on terminals. The shown targets $t_-,t_+$ and hole endpoints
    $h_1,h_2$ are retained; the terminal-free branch has zero flow, and each
    surviving path becomes one segment with the same coefficient. This is
    the tree compression used in constructing the remainder for
    Lemma~\ref{lem:backend-reduction}; the pictured contribution need not be
    closed. Vertices shown here are copied tree nodes. Physical hubs, null
    connectors, selected chords, and the contributions of every other
    frozen and raw forest join the same child graph $J$.}
  \label{fig:frozen-compression}
\end{figure}


\begin{samepage}
\begin{lemma}[One-level quality reduction]
In every state of an epoch, suppose the active logical path
assignments use selected chords and have at most $p_*$ hops.
\label{lem:backend-reduction}
Let $L_{\mathrm{best}}$ be the largest local correction score
\eqref{back:local-score}, or zero if there is none.  Then, for every
nonzero circulation $\Delta$ on the current nonloops,
\begin{equation}
\label{back:reduction-bound}
 \frac{|\gamma^{\mathsf T}\Delta|}{\|\ell\Delta\|_1}
 \le 2Z_*A_*L_{\mathrm{best}}+12Z_*A_*\OPT(J)
 \le b^2\max\{L_{\mathrm{best}},\OPT(J)\}.
\end{equation}
This conclusion is simultaneous for all current circulations and
requires no independence between the state and the verified covers
or path assignments.
\end{lemma}
\end{samepage}

\begin{proof}
Use the coefficients and raw vector of
Lemma~\ref{back:frozen-decomposition}, and form on the child edges
\begin{equation}
\label{back:child-vector}
                  \Xi_0=\sum_f\alpha_fJ_f-\widehat z_L+\Delta_n,
\end{equation}
with $\Delta_n$ carried by the hub-to-hub new edges.  Its physical
projection is
\[
                \pi\Xi_0=\Delta-\sum_f\alpha_f\pi B_f,
\]
a physical circulation.  Since $W_f\le2w_f$, the norm estimates give
\begin{align}
 \|\ell^J\Xi_0\|_1
 &\le 2Z_*A_*\|\ell\Delta_s\|_1
            +A_*\|\ell\Delta_s\|_1+\|\ell\Delta_n\|_1\notag\\
 &\le4Z_*A_*\|\ell\Delta\|_1.
                                    \label{back:child-norm-before}
\end{align}
The vector $\Xi_0$ need not balance at each compressed copy of a
physical vertex.  Route its imbalance there along the null edge to
that vertex's hub.  The hub then balances because the sum of
divergences over all copies with its physical image is zero.
The total absolute imbalance at the compressed vertices is at most
twice the sum of absolute coefficients of $\Xi_0$, and hence at
most twice its length norm, since every child edge has length at
least one.  The resulting child circulation $\Xi$ therefore satisfies
\begin{equation}
\label{back:child-norm-after}
 \pi\Xi=\pi\Xi_0,\qquad
              \|\ell^J\Xi\|_1\le12Z_*A_*\|\ell\Delta\|_1.
\end{equation}
Its gradient equals that of its projection because all connectors
have zero gradient.

For the local part, \eqref{back:local-score} and
\eqref{back:coefficient-bound} give
\[
 \left|\sum_f\alpha_f\gamma^{\mathsf T}\pi B_f\right|
 \le L_{\mathrm{best}}Z_*\sum_f|\alpha_f|W_f
 \le2Z_*A_*L_{\mathrm{best}}\|\ell\Delta\|_1.
\]
The child circulation contributes at most its norm times $\OPT(J)$,
with the same conclusion if it is zero.  Adding these inequalities
proves the first bound in \eqref{back:reduction-bound}.  The second
follows from $14Z_*A_*\le20Z_*A_*\le b^2$.

Every distribution, coupling, coefficient, and comparison vector
in this proof is existential.  The algorithm constructs only
integral graphs, terminal sets, path assignments, and rational
scores.  Once the covers and assignments have their verified
properties, the argument applies to any state and any circulation
chosen afterwards.  This proves the asserted adaptive scope.
\end{proof}

\begin{corollary}[Quality through the recursion]
\label{back:quality-induction}
Assume the path assignments and public forest lifting of
Section~\ref{sec:backend-updates}.  The backend at level $i$ returns
a certified score at least
\begin{equation}
\label{back:level-score}
                    \frac12 b^{-2(d-i)}\OPT_i.
\end{equation}
At the outer level this is at least $M^{-10/a}\OPT$.
\end{corollary}

\begin{proof}
At the bottom use Lemma~\ref{back:base}.  At a higher level compare
the local correction scores, exact current-loop scores, and the
recursive child score.  Lemma~\ref{update:flat-lift} preserves the
child gradient under lifting and does not increase raw actual length,
so the child's certified score remains valid.  If its approximation
factor is $\eta\le1$, the maximum of its score and
$L_{\mathrm{best}}$ is at least
$\eta\max\{L_{\mathrm{best}},\OPT(J)\}$.  Apply
Lemma~\ref{lem:backend-reduction} to lose at most $b^2$ at this
level.  Exact loop candidates give the same conclusion for the full
graph.  This proves \eqref{back:level-score} by induction.

Since $d=2a-4$, the outer factor is
$\tfrac12M^{-(4a-8)/a^2}$, which exceeds $M^{-10/a}$ in the checked
range.  A zero optimum does not require a positive answer.  The
representation size, explicit forest logs, and total operation cost
are established next; the quality argument itself has incurred only
one recursive child call at each level.
\end{proof}
